\phantomsection
\chapter*{Önsöz}
\markboth{Önsöz}{Önsöz}
\addcontentsline{toc}{chapter}{Önsöz}

Bu kitap, PDR 209 Temel İstatistik dersinin öğrenme çıktıları doğrultusunda
hazırlanmıştır. Temel amaç, öğrencinin yalnızca hesaplama yapması değil;
araştırma sorusunu istatistiksel bir probleme dönüştürmesi, uygun yöntemi
seçmesi, belirsizliği yorumlaması ve bulguyu PDR bağlamında sorumlu biçimde
raporlamasıdır.

Konular önkoşul ilişkileri gözetilerek sıralanmıştır. Önce veri üretimi,
ölçme ve örnekleme ele alınır; ardından betimsel istatistik, normal dağılım,
tahmin ve güven aralıkları kurulur. Hipotez testleri bu temel üzerinde
geliştirilir; t testleri, ANOVA, korelasyon, regresyon, kategorik ve
parametrik olmayan yöntemler araştırma sorularıyla ilişkilendirilir. Son
bölümde bütün süreç tek bir veri analizi laboratuvarında birleştirilir.

Her bölüm dört aşamalı bir öğrenme döngüsü izler: hedefi görme, kavramı
anlama, kısa bir probleme uygulama ve sonucu bağlam içinde yorumlama. Bölüm
sonu soruları yalnız formül hatırlamayı değil; yöntem seçme, hata fark etme,
çıktı okuma ve etik karar verme becerilerini de yoklar.

Kitapla çalışırken önce bölümün öğrenme hedeflerini okuyunuz; örneklerdeki
hesaplamaları kâğıt üzerinde tahmin etmeye çalışınız; ardından yazılım
çıktısıyla karşılaştırınız. Sonucun yönünü ve yaklaşık büyüklüğünü yazılımı
çalıştırmadan öngörebilmek, istatistiksel sezginin önemli bir göstergesidir.
Her analizde şu dört soruyu tekrar ediniz: Araştırma sorusu nedir? Veri nasıl
üretilmiştir? Yöntemin koşulları sağlanıyor mu? Sonuç PDR uygulamasında ne
anlama gelmektedir?

Kitaba eşlik eden dijital uygulamalara
\url{https://ibrahimguney.github.io/istatistik-lab/} adresinden erişilebilir.
Bu uygulamalar, sonuç üretmekten çok analiz kararlarını görünür kılmak ve
hesaplamaları yeniden üretilebilir hâle getirmek amacıyla kullanılmalıdır.

\begin{flushright}
Prof. Dr. İbrahim Güney\\
İstanbul, 2026
\end{flushright}